\section{Definición del Problema}

En un laboratorio de Control, el objetivo central del estudiante es dominar el modelado matemático del sistema, simularlo con herramientas científicas (como Matlab/Simulink) y contrastar esos resultados teóricos con datos experimentales obtenidos de la planta física. \\
Estos procesos pedagógicos son obstaculizados cuando el estudiante debe invertir su tiempo de clase en tareas técnicas secundarias, como configurar sensores, programar microcontroladores o gestionar cables; actividades que no aportan valor al aprendizaje de conceptos fundamentales de control.\\\\

La solución propuesta \textbf{una interfaz web de adquisición de datos} elimina estas barreras. Con solo acceder a una URL desde cualquier dispositivo (computadora o celular), el estudiante puede iniciar el experimento, visualizar las variables en tiempo real y descargar los datos en formato compatible (CSV, Excel, txt) para su análisis, sin interactuar con hardware, instalar software o depurar código.\\
Esta simplicidad contrasta con alternativas existentes. Plataformas como Arduino aunque exigen mínimos conocimientos de electrónica y programación, quitan tiempo relevante durante clase.\\
Por otro lado, soluciones profesionales como NI ELVIS o MATLAB Data Acquisition Toolbox, aunque potentes, tienen costos prohibitivos y complejidades técnicas que necesitan ser atendidas previamente.\\\\

Además de la accesibilidad del sistema propuesto, destaca por su escalabilidad y bajo costo. Por menos de \$500 MXN se consigue una Raspberry Pi capaz de implementar una infraestructura como la planteada. Su arquitectura modular gracias al MVC permite mejores futuras sin modificar la base del código.\\
Esto no solo optimiza recursos institucionales, sino que asegura que las actualizaciones se alineen con las necesidades evolutivas de la enseñanza.
